\section{Discussion}
\label{sec:discussion}

Our dual-track design separates reproducible mechanism tests from deployment realism. Public benchmarks allow independent verification of routing and retrieval, whereas CNAS deployment shows that the full pipeline operates under latency and HITL constraints in an accredited laboratory. Public corpora cannot capture CNAS-specific numeric and seal checks; conversely, the Tier-1 proxy does not recover deployment accuracy on the temporal hold-out.

The priority cascade is best interpreted conditionally. On rule-routed items, engine conflict remains substantial, yet lenient unsupported positives fall to zero under cascade while OR-ensemble still reports small rates. Strict counting shows that the rule-engine false positives affect both designs equally, so the cascade is best justified by generative-risk isolation, rule traceability, and reduced HITL load rather than by minimizing every false positive. The safety benefit is conditional: it trades strict FP equivalence for traceability and generative risk isolation on routed items. On full ContractNLI, OR-ensemble can achieve lower overall lenient hallucination by fusing RAG on semantic-only items, whereas CNAS deployment benefits from gating through lower HITL load and lower lenient hallucination than ensemble or RAG-only baselines. Self-RAG and ReAct were evaluated on C3PA only (Appendix~\ref{app:tier2_transfer}); CNAS comparisons are limited to Rule, RAG, Ensemble, and \HRAG{}. \HRAG{}'s gating reduces average LLM calls to 1.3 per report by routing 95\% of checklist items through the rule engine, which explains its lower latency and HITL load despite comparable retrieval quality.

Relative to hybrid ACC systems that combine semantic alignment with executable rule interpretation \cite{zhang2016acc,zheng2022acc} and regulation-aware datasets \cite{c3pa2024,koreeda2021contractnli}, we emphasize explicit priority routing with stratified evaluation and deployment-side HITL governance. Compared with retrieve--rerank--generate pipelines and self-critique RAG variants \cite{glass2022re2g,asai2023selfrag}, our priority cascade targets checklist-level auditing with auditable rule traces rather than open-domain answer fusion.

Several limitations remain: idealized rule coverage, a small CNAS corpus (87 reports), English public benchmarks versus Chinese deployment reports, interpretive rather than bound-based analysis, sensitivity between lenient and strict hallucination metrics (\ref{app:hallucination}), and brittleness under regulatory version drift (Tier-2b, 65.3\%). In production, rules absorb most checklist items, reducing average LLM calls to 1.3 per report and concentrating human review on low-confidence outputs while preserving evidence traces---but only while the encoded rule set matches the current standard edition.
